\documentclass[11pt]{article}

\usepackage[margin=1in]{geometry}
\setlength{\emergencystretch}{2em}
\usepackage{amsmath,amssymb,amsthm}
\usepackage{hyperref}

% ---------------------------------------------------------------
% Placeholder macro: everything that must be filled in by hand.
% Search for "TODO" before sending.
% ---------------------------------------------------------------
\newcommand{\todo}[1]{\textbf{[TODO: #1]}}

% ---------------------------------------------------------------
% Notation, as in Tschinkel--Yang--Zhang
% ---------------------------------------------------------------
\newcommand{\BC}{\mathcal{BC}}   % combinatorial Burnside group  BC_n(G)
\newcommand{\SC}{\mathcal{SC}}   % combinatorial symbols group    SC_n(G)
\newcommand{\B}{\mathcal{B}}     % symbols group of Kontsevich--Pestun--Tschinkel B_n(H)
\newcommand{\Sym}{\mathfrak{S}}  % symmetric group  S_m
\newcommand{\Alt}{\mathfrak{A}}  % alternating group
\newcommand{\Z}{\mathbb{Z}}
\newcommand{\Q}{\mathbb{Q}}
\newcommand{\Cx}{\mathbb{C}^{\times}}
\newcommand{\dual}[1]{#1^{\vee}}
\newcommand{\Cen}[2]{Z_{#1}(#2)}  % centralizer  Z_G(H)
\newcommand{\Nor}[2]{N_{#1}(#2)}  % normalizer   N_G(H)

\title{Computing combinatorial Burnside groups $\BC_n(G)$ in Julia/OSCAR\\[2pt]
       \large A working note}
\author{\todo{author name(s)}}
\date{\today}

\begin{document}
\maketitle

\section{Introduction}

This note reports on a computer implementation, in Julia via the OSCAR
system~\cite{OSCAR} (which calls GAP~\cite{GAP} for the group-theoretic part), of the
combinatorial Burnside groups $\BC_n(G)$ of Tschinkel--Yang--Zhang~\cite{TYZ}, for finite
groups $G$ given as permutation or matrix groups. The goal is to reproduce the published
computations of~\cite{TYZ} and, where feasible, to extend them to further groups and
values of~$n$. We use the notation of~\cite{TYZ} throughout: symbols
$(H,Y,\beta)\in\SC_n(G)$, the groups $\B_n(H)$, and the relations (O), (C), (V), (B2), and
(B2$'$).

\section{Method}

Let $G$ be a finite group and $n\ge 1$. By Theorem~5.2 of~\cite{TYZ} together with
Lemma~5.1,
\begin{equation}\label{eq:decomp}
  \BC_n(G)\;\cong\;\bigoplus_{[H,Y]}\B_n([H,Y]),
  \qquad
  \B_n([H,Y])\;\cong\;\B_n(H)\big/(C_{(H,Y)}),
\end{equation}
where the sum runs over $G$-conjugacy classes of pairs $(H,Y)$ with $H\subseteq G$ abelian
(and nontrivial) and $H\subseteq Y\subseteq \Cen{G}{H}$, and $(C_{(H,Y)})$ is the relation
$\beta=\beta^g$ for $g\in \Nor{G}{H}\cap \Nor{G}{Y}$. The algorithm computes each summand
explicitly and adds them up.

\paragraph{Step 1: classes $[H,Y]$ (GAP).}
We enumerate representatives $H$ of the conjugacy classes of nontrivial abelian subgroups
of~$G$. For each $H$ we enumerate the subgroups $Y$ with $H\subseteq Y\subseteq \Cen{G}{H}$
and take their orbits under $\Nor{G}{H}$; each orbit is one class $[H,Y]$, with stabilizer
$\Nor{G}{H}\cap \Nor{G}{Y}$. Writing $H=\langle g_1\rangle\times\dots\times\langle g_k\rangle$,
the action of the stabilizer on $H$, and hence on $\dual{H}$, is recorded as integer
matrices (exponents of $g_j^{\,g}$ in terms of the $g_i$).

\paragraph{Step 2: the group $\B_n(H)$.}
We write $\dual{H}$ additively. A generator of $\B_n(H)$ is an unordered $n$-tuple
$\beta=(b_1,\dots,b_n)$ of characters, \emph{zeros allowed}, that generate $\dual{H}$; this
is the zero-padded form used in the proof of Lemma~5.1 of~\cite{TYZ}, where a symbol
$(H,Y,\beta)$ of length $r\le n$ corresponds to $\beta$ padded by $n-r$ zeros. This
imposes (O). The relations are the following, for each pair of positions $i<j$ with
$a=b_i$, $b=b_j$:
\begin{itemize}
  \item (B): if $a\ne b$ and both are nonzero, $\beta=\beta_1+\beta_2$ with
        $\beta_1=(\dots,a-b,\dots,b,\dots)$ and $\beta_2=(\dots,a,\dots,b-a,\dots)$;
  \item if $a=b\neq 0$, $(b,b,\dots)=(0,b,\dots)$ (the case $b_1=b_2$ of (B2$'$));
  \item (V): if $a+b=0$ with $a,b\ne 0$, then $\beta=0$;
  \item $(C_{(H,Y)})$: $\beta=\beta^g$ for $g$ in the (image of the) stabilizer.
\end{itemize}
Note that (B) is applied to nonzero pairs only: applied literally to $b_1=b_2=0$ it would force
$\beta=2\beta$. \todo{confirm with Prof.\ Tschinkel that this is the intended convention
for $\B_n$ as in the proof of Lemma~5.1.}

\paragraph{Step 3: the abelian group.}
The relations form an integer matrix whose cokernel is $\B_n(H)/(C_{(H,Y)})$. We first
eliminate generators using relations with a coefficient $\pm1$, and then compute the Smith
normal form of the (small) remaining matrix over~$\Z$. This yields the invariant factors,
which we convert to the primary decomposition (a product of groups $\Z/p^e$ and copies of $\Z$) used in~\cite{TYZ}.
Summands are cached by the isomorphism type of $H$ and the image of the stabilizer in
$\mathrm{Aut}(H)$.

\paragraph{Step 4: assembly.}
Summing the groups over all classes $[H,Y]$ as in~\eqref{eq:decomp} gives $\BC_n(G)$.

\paragraph{Independent check.}
As a check that does not rely on Theorem~5.2, we also implemented $\BC_n(G)$ directly from
the definition: generators $(H,Y,\beta)$ over \emph{all} abelian $H$ and \emph{all} $Y$ (not up
to conjugacy), with the relations (O), (C), (V) and (B2), including the term
$\Theta_2=(\bar H,Y,\bar\beta)$, $\bar H=\ker(b_1-b_2)$, $\bar\beta=\beta|_{\bar H}$. This is
only practical for small groups. Groups for which the two implementations were compared:
\todo{list of groups and values of $n$}.

\section{Results}

Table~\ref{tab:results} compares the values in~\cite{TYZ} with our computed values.
Groups are given as in~\cite{TYZ}; values are written in the primary decomposition
notation of that paper.

\begin{table}[ht]
\centering
\begin{tabular}{|c|c|l|l|c|}
\hline
$G$ & $n$ & Value in \cite{TYZ} & Computed value & Agree? \\
\hline
\todo{G} & \todo{n} & \todo{paper} & \todo{ours} & \todo{y/n} \\
\todo{G} & \todo{n} & \todo{paper} & \todo{ours} & \todo{y/n} \\
\todo{G} & \todo{n} & \todo{paper} & \todo{ours} & \todo{y/n} \\
\hline
\end{tabular}
\caption{Published versus computed values of $\BC_n(G)$.
\todo{add/remove rows; note which entries, if any, disagree and why.}}
\label{tab:results}
\end{table}

\todo{Optional: values computed for groups or $n$ not covered in~\cite{TYZ}, clearly
separated from the table above.}

\section{Scope and limitations}

The computation has two potential bottlenecks. The first is the enumeration in GAP of
conjugacy classes of subgroups of~$G$ (Step~1). The second is the size of $\B_n(H)$,
whose number of generators is $\binom{|H|+n-1}{n}$ before the relations are imposed
(Step~2). The largest groups and values of $n$ we have handled so far are:
\todo{largest $|G|$ / groups tractable, with timings}; and the point at which
each bottleneck became prohibitive was \todo{observed thresholds}.
The independent direct implementation is limited to small groups (\todo{observed range}).

We have not implemented the restriction maps or the ring structure on
$\BC_*(G)$ from~\cite{TYZ}, nor the map $\mathrm{Burn}_n(G)\to\BC_n(G)$ used
in~\S6.5 there.

\section{Code availability}

\todo{repository or gist link, and the versions of Julia and OSCAR used.}

\begin{thebibliography}{9}
\bibitem{TYZ} Yuri Tschinkel, Kaiqi Yang, and Zhijia Zhang,
  \emph{Combinatorial Burnside groups}, arXiv:2112.12801,
  \url{https://arxiv.org/abs/2112.12801}.
\bibitem{OSCAR} The OSCAR team, \emph{OSCAR -- Open Source Computer Algebra Research system},
  \url{https://www.oscar-system.org}.
\bibitem{GAP} The GAP Group, \emph{GAP -- Groups, Algorithms, and Programming},
  \url{https://www.gap-system.org}.
\end{thebibliography}

\end{document}
